\documentclass[11pt]{amsart}
\usepackage{calc,amssymb,amsmath,ulem,stmaryrd,fullpage}
\usepackage{enumerate}
\usepackage{alltt}
\RequirePackage[dvipsnames,usenames]{color}
\let\mffont=\sf
\normalem
%\input{kmacros3.sty}
\input{xy}
\xyoption{all}
\usepackage{hyperref}
\usepackage{graphicx}
\usepackage[all,cmtip]{xy}
\usepackage{verbatim}
\usepackage[left=0.95in,top=0.95in,right=0.95in,bottom=0.95in]{geometry}
\newcommand{\tauCohomology}{T}
\newcommand{\FCohomology}{S}


\theoremstyle{theorem}
%\newtheorem*{mainthm*}{Main Theorem}

\newcommand{\note}[1]{\marginpar{\sffamily\tiny #1}}

\def\todo#1{\textcolor{Mahogany}%
{\sffamily\footnotesize\newline{\color{Mahogany}\fbox{\parbox{\textwidth-15pt}{#1}}}\newline}}

\renewcommand{\O}{\mathcal O}
\newcommand{\ie}{{\itshape i.e.} }
\newcommand{\roundup}[1]{\lceil #1 \rceil}
\newcommand{\rounddown}[1]{\lfloor #1 \rfloor}
\renewcommand{\to}[1][]{\xrightarrow{\ #1\ }}
\newcommand{\onto}[1][]{\protect{\xrightarrow{\ #1\ }\hspace{-0.8em}\rightarrow}}
\newcommand{\into}[1][]{\lhook \joinrel \xrightarrow{\ #1\ }}
%\newcommand{DBDual}[1]{{\underline{\omega}_{#1}}}
\newcommand{\DBDual}[1]{{{\uuline \omega}{}^{\mydot}_{#1}}}
\newcommand{\Tt}{{\mathfrak{T}}}
%\newcommand{\bn}{{\mathfrak{n}} }
\newcommand{\sol}[1]{ \vskip 6pt{\color{blue} {\bf Solution:} #1} \vskip 6pt}

\newcommand{\bR}{{\mathbb{R}}}
\newcommand{\va}{{\vec{a}}}
\newcommand{\vb}{{\vec{b}}}
\newcommand{\vzero}{{\vec{0}}}
\newcommand{\vi}{{\vec{i}}}
\newcommand{\vj}{{\vec{j}}}
\newcommand{\vk}{{\vec{k}}}
\newcommand{\vh}{{\vec{h}}}
\newcommand{\vr}{{\vec{r}}}
\newcommand{\vu}{{\vec{u}}}
\newcommand{\vv}{{\vec{v}}}
\newcommand{\vw}{{\vec{w}}}
\newcommand{\vx}{{\vec{x}}}
\newcommand{\vy}{{\vec{y}}}
\newcommand{\vz}{{\vec{z}}}
\newcommand{\vT}{{\vec{T}}}
\newcommand{\vN}{{\vec{N}}}
\newcommand{\vF}{{\vec{F}}}
\newcommand{\vG}{{\vec{G}}}
\newcommand{\bF}{\mathbb{F}}
\newcommand{\bQ}{\mathbb{Q}}
\newcommand{\bZ}{\mathbb{Z}}
\newcommand{\bC}{\mathbb{C}}
\newcommand{\Gal}{\textnormal{Gal}}
\newcommand{\Aut}{\textnormal{Aut}}
\newcommand{\Emb}{\textnormal{Emb}}
\newcommand{\Tr}{\textnormal{Tr}}
\newcommand{\N}{\textnormal{N}}


\begin{document}
\title{HW \#3 -- Math 6320\\Spring 2015}
\author{Due: Tuesday February 17th}
\maketitle

\begin{enumerate}
\item Exhibit an \emph{explicit} isomorphism between the splitting fields of $x^3 - x + 1$ and $x^3 - x - 1$ over $\bF_3$. 
\item Explicitly show that $\bQ(\sqrt{2}, \sqrt{3}, \sqrt{5})$ is a simple extension of $\bQ$ by finding a generator.
\item Show that $\bF_p(x,y) / \bF_p(x^p, y^p)$ is \emph{not} a simple extension.   
\item Is $\bQ(2^{1/3})$ a subfield of a cyclotomic extension of $\bQ$?
\item Let $K/F$ be any finite extension, let $L$ be a Galois extension of $F$ containing $K$ and let $H \leq \Gal(L/F)$ be the subgroup corresponding to $K$.  For any $\alpha \in K$ define the \emph{norm} of $\alpha$ to be
    \[
    {\N}_{K/F}(\alpha) = \prod_{\overline{\sigma} \in G/H} \overline{\sigma}(\alpha).
    \]
    Note here that $G/H$ is just the set of cosets of $H$.  Likewise we define the \emph{trace} of $\alpha$ to be 
    \[
    {\Tr}_{K/F}(\alpha) = \sum_{\overline{\sigma} \in G/H} \overline{\sigma}(\alpha).
    \]
    In other words, the norm is a product of Galois conjugates and the trace is a sum of Galois conjugates (if $K/F$ is Galois, it is a product/sum over all Galois conjugates).  Of course,
    \begin{itemize}
    \item[(a)]  Show that $\N_{K/F}(\alpha) \in F$ and $\Tr_{K/F}(\alpha) \in F$.  
    \item[(b)]  Prove that the norm is multiplicative ($\N_{K/F}(\alpha \beta) = \N_{K/F}(\alpha) \N_{K/F}(\beta)$) and that the trace is additive ($\Tr_{K/F}(\alpha \beta) = \Tr_{K/F}(\alpha) + \Tr_{K/F}(\beta)$)
    \item[(c)]  If $K = L = F(\sqrt{c})$ is a degree two extension, find formulas for $\N(a + b\sqrt{c})$ and $\Tr(a + b \sqrt{c})$.
    \end{itemize}
\item Let $K/F$ be a finite separable extension.  Then each element $\alpha \in K$ gives us an $F$-linear map $(\times \alpha) : K \to K$.  We can of course pick a basis for $K/F$ and so consider $(\times \alpha)$ as a matrix $m_{\alpha}$.  Then consider maps $t : K \to F$ and $n : K \to F$ defined by
    \[
    t(\alpha) = {\textnormal{trace}}(m_{\alpha}), n(\alpha) = {\textnormal{det}}(m_{\alpha})
    \]
    where here trace just means sum up the diagonal.  Is it true that $t = \Tr_{K/F}$ and $n = \N_{K/F}$?
\item Fix notation as in the last two problems.  If $\alpha \in F$, show that $\N(\alpha) = \alpha^n$ and $\Tr(\alpha) = n\alpha$ where $n = [K:F]$.
\item Fix notation as in the last three problems.  Let $F = \bQ, K = \bQ(2^{1/3})$ and $\alpha = 2^{1/3}$.  Compute ${\N}_{K/F}(2^{1/3})$ and ${\Tr}_{K/F}(2^{1/3})$.
\end{enumerate}

\end{document}

